% fig4m.tex: Fig. 4, variant 4M (4L as a measurement engineer would read it: saturation marked, the known
% systematic part shown, intervals that hold for all cells at once). One column, 252 pt wide.
%   \input{figures/tikz_v2/fig4m.tex}   inside \begin{figure}[t]\centering ... \end{figure}
% Preamble: \usepackage{pgfplots}\pgfplotsset{compat=1.18}% icorr_style.tex: shared colours and line styles of the amplitude-free paper's figures
% (paper/notes/story_concept.md section 7a, scheme C). \input once in the preamble, after
% \usepackage{pgfplots}. Every figure uses these names; no figure defines its own colours.
%
% Roles:
%   PsyPref (all variants) ............ icBlue, solid
%   best published comparator ......... icOrange, solid, hollow marker: square CMA-ES, triangle Arens' GP
%   GP-BALD (same prior) .............. icGrey, dashed
%   CMA-ES on Schaefer's family ....... icDark, dash-dot
%   generic assistance (50 %) ......... icGeneric, solid thin
%   best K-weight profile (ceiling, floor) icContext, dotted
%   K = 8 / 16 / 24 (or k = 24 / 40 / 150) icBlueLight / icBlue / icBlueDark
%   settings as markers ............... circle 3-35 Nm, square fixed 12 Nm, diamond 3-12 Nm
%   robustness: filled = fixed yardstick (post hoc), hollow = registered scoring
\definecolor{icBlue}{HTML}{1845FB}
\definecolor{icBlueLight}{HTML}{8FA4FD}
\definecolor{icBlueDark}{HTML}{0F2C9E}
\definecolor{icOrange}{HTML}{F89C20}
\definecolor{icOrangeText}{HTML}{B86E00}% text in the comparator's hue (contrast 4.0 on white)
\definecolor{icGrey}{HTML}{7A7A7A}
\definecolor{icDark}{HTML}{333333}
\definecolor{icGeneric}{HTML}{9C9CA1}
\definecolor{icContext}{HTML}{BBBBBB}
\definecolor{icRule}{HTML}{DDDDDD}
\definecolor{icInk}{HTML}{222222}
\definecolor{icMuted}{HTML}{666666}
% Fig. 1 only: the pair shown to the user (author's choice, 2026-10-29: distinct hues read
% stronger than two blues; gold was dropped because it equals the comparator's orange)
\definecolor{icTintBlue}{HTML}{F1F3F8}% Fig. 1: the loop, PsyPref's and the user's box (author's choice)
\definecolor{icPairA}{HTML}{BB5566}
\definecolor{icPairB}{HTML}{33BBEE}
\tikzset{
  ic psypref/.style={icBlue, line width=1.5pt, line cap=round, line join=round},
  ic best/.style={icOrange, line width=1.2pt, line cap=round, line join=round},
  ic gpbald/.style={icGrey, line width=0.9pt, dash pattern=on 3.2pt off 1.6pt},
  ic sch/.style={icDark, line width=0.9pt, dash pattern=on 3.6pt off 1.4pt on 0.9pt off 1.4pt},
  ic generic/.style={icGeneric, line width=0.6pt},
  ic context/.style={icContext, line width=0.8pt, dash pattern=on 0.9pt off 1.5pt},
  ic rule/.style={icRule, line width=0.5pt},
  ic label/.style={inner sep=0.6pt, font=\scriptsize, align=left},
  ic note/.style={inner sep=0.8pt, font=\scriptsize, icInk, fill=white, fill opacity=0.85,
                  text opacity=1, align=right},
  ic band/.style={fill=#1, fill opacity=0.15, draw=none},
}
\pgfplotsset{
  ic axis/.style={
    axis line style={icMuted, line width=0.5pt},
    tick style={icMuted, line width=0.5pt},
    tick label style={font=\scriptsize, icMuted},
    label style={font=\footnotesize, icInk},
    title style={font=\footnotesize, icInk},
    axis x line*=bottom, axis y line*=left,
    tick align=outside, major tick length=2pt,
    clip=false,
  },
}

% Data: \icorrtikzdata (default figures/tikz_v2/data/), written by scripts/paper_v2_fig4m.py (amp-free
% repository): fig4m_rows.csv (one line per group header or row: y in row units; per setting the light and
% the blue five values, nan where the cell has the other colour or does not exist; the cap marker; the
% label), fig4m_meta.csv (axis limits and the layout in pt).
% Forest plot: PsyPref minus the strongest optimizer of that row, picked again in every bootstrap draw,
% paired by user, chosen over generic assistance after 24 comparisons, in points, each user rated with their
% own noise level and lapse rate at the session's end. Circle 3-35 Nm (0.25 rows above the row's centre),
% square fixed 12 Nm (0.25 below). Thick bar: 95 % interval of that row and setting. Thin bar: 95 % for the 73
% cells jointly. Blue with its value where the thin bar reaches 0, else icBlueLight. Open circle with "both
% at 95 %": PsyPref and the comparator at the measure's maximum in every draw. "(n)": the smallest lead of n
% conditions per setting; dagger: PsyPref with tie handling (n = 160). The first and the last row are
% bold, no header; the second row is the main condition against Fig. 2's comparator. Row labels give a
% deviation's size where one number suffices.
\providecommand{\icorrtikzdata}{figures/tikz_v2/data/}%
\makeatletter
\providecommand{\icorrload}[2]{% \icorrload{name}{file.csv}: every cell as \icv{name}{row}{column}%
  \pgfplotstableread[col sep=comma]{\icorrtikzdata #2}\ic@table
  \pgfplotstableforeachcolumn\ic@table\as\ic@col{%
    \pgfplotstableforeachcolumnelement{\ic@col}\of\ic@table\as\ic@cell{%
      \expandafter\global\expandafter\let\csname ic@#1@\pgfplotstablerow @\ic@col\endcsname
        \ic@cell}}}%
\providecommand{\icv}[3]{\csname ic@#1@#2@#3\endcsname}%
\def\fmlabel#1#2{\expandafter\def\csname fm@lab@#1\endcsname{#2}}%
\def\fmlab#1{\csname fm@lab@#1\endcsname}%
\makeatother
% the row labels, by key (fig4m_rows.csv's column key); condition values from fig4m_meta.csv
\fmlabel{MAIN}{Main condition (Fig.~2)}%
\fmlabel{G1}{How they answer}%
\fmlabel{R1}{Noise level $\sigma$ = \pgfmathprintnumber{\icv{fmm}{0}{s_r1}}\,Nm}%
\fmlabel{R2}{Noise level $\sigma$ = \pgfmathprintnumber{\icv{fmm}{0}{s_r2}}, \pgfmathprintnumber{\icv{fmm}{0}{s_r3}}\,Nm}%
\fmlabel{AL}{Lapse rate $\lambda$ = \pgfmathprintnumber{\icv{fmm}{0}{lapse_a2}}, \pgfmathprintnumber{\icv{fmm}{0}{lapse_r3}}}%
\fmlabel{A1}{$\sigma$ differs per user (log SD \icv{fmm}{0}{a1_logsd})}%
\fmlabel{A6}{Other answer-curve shapes}%
\fmlabel{A3}{Device noise \icv{fmm}{0}{a3_noise}\,Nm}%
\fmlabel{A4}{Bias toward first profile, \icv{fmm}{0}{a4_bias}\,$\sigma$}%
\fmlabel{A5T}{``No preference'' below \icv{fmm}{0}{a5_nopref}\,JND}%
\fmlabel{G2}{How they change}%
\fmlabel{A7}{Fatigue: $\lambda$ from \icv{fmm}{0}{a7_from} to \icv{fmm}{0}{a7_to}}%
\fmlabel{A8}{$\sigma$ \icv{fmm}{0}{a8_from} to \icv{fmm}{0}{a8_to}\,Nm, and reverse}%
\fmlabel{B4}{Preferred profile drifts \icv{fmm}{0}{b4_a}, \icv{fmm}{0}{b4_b}\,JND}%
\fmlabel{B3}{Familiarity}%
\fmlabel{G3}{What they judge and prefer}%
\fmlabel{B1}{Weber's law}%
\fmlabel{B2}{Over-assistance counts double}%
\fmlabel{B5}{Judging single features}%
\fmlabel{B67}{Gait phases weighted unequally}%
\fmlabel{B8}{No single preferred profile}%
\fmlabel{R4}{Negative power disliked}%
\fmlabel{G4}{How they differ from the prior}%
\fmlabel{R6}{Users from other datasets}%
\fmlabel{PM}{Peaks unlike the prior}%
\fmlabel{PS}{Flexion peak 6, 12\,\%GC later}%
\fmlabel{SO}{Other shapes unlike the prior}%
\fmlabel{SC}{Shapes from the 4--6-param. family}%
\fmlabel{C4T}{Pessimistic user (\icv{fmm}{0}{ndev_word} at once)}%
\begin{tikzpicture}[x=1pt, y=1pt]
\pgfkeys{/pgf/number format/assume math mode=true}% numbers in the text font (Times)
\icorrload{fmm}{fig4m_meta.csv}%
\icorrload{fmr}{fig4m_rows.csv}%
\pgfmathtruncatemacro\fmN{\icv{fmm}{0}{n_rows}-1}%
\pgfmathsetmacro\fmL{\icv{fmm}{0}{lab_w}}\pgfmathsetmacro\fmR{\icv{fmm}{0}{ax_r}}%
\pgfmathsetmacro\fmP{0.85*\icv{fmm}{0}{pitch}}\pgfmathsetmacro\fmB{\icv{fmm}{0}{bottom}}%
\pgfmathsetmacro\fmoff{\icv{fmm}{0}{off}}\pgfmathsetmacro\fmmr{\icv{fmm}{0}{mark_r}}%
\pgfmathsetmacro\fmms{\icv{fmm}{0}{mark_s}}\pgfmathsetmacro\fmvgapx{\icv{fmm}{0}{val_gapx}}%
\pgfmathsetmacro\fmxmin{\icv{fmm}{0}{xmin}}\pgfmathsetmacro\fmxmax{\icv{fmm}{0}{xmax}}%
\pgfmathsetmacro\fmymin{\icv{fmm}{0}{ymin}}\pgfmathsetmacro\fmymax{\icv{fmm}{0}{ymax}}%
\pgfmathsetmacro\fmscale{\icv{fmm}{0}{scale}}%
\pgfmathsetmacro\fmH{(\fmymax-\fmymin)*\fmP}\pgfmathsetmacro\fmW{\fmR-\fmL}%
\pgfmathsetmacro\fmylab{\icv{fmm}{0}{ylabels}}\pgfmathsetmacro\fmlinetop{\icv{fmm}{0}{linetop}}%
\def\fmlight{icBlueLight}% a lead whose thin bar lies above 0: the lighter PsyPref
\pgfplotsset{
  fm mark w/.style={mark=*, mark size=\fmmr pt},
  fm mark f/.style={mark=square*, mark size=\fmms pt},
  fm marks/.style={only marks, mark options={fill=#1, draw=#1, line width=0.3pt}},
  % error bars only: no line through the points and no marks
  fm bar/.style 2 args={draw=none, mark=none, error bars/.cd, x dir=both, x explicit, error mark=none,
    error bar style={draw=#1, line width=#2, line cap=butt}, /pgfplots/.cd},
}
\begin{axis}[ic axis, scale only axis, at={(\fmL pt,\fmB pt)}, anchor=south west,
  width=\fmW pt, height=\fmH pt,
  xmin=\fmxmin, xmax=\fmxmax, ymin=\fmymin, ymax=\fmymax, y dir=reverse,
  axis y line=none, ytick=\empty,
  xtick={-10,0,10,20,30},
  xticklabels={\textminus10, 0, +10, +20, +30},
  xticklabel style={inner sep=1pt},
]
% no lead (0): the strongest published optimizer
\draw[icMuted, line width=0.5pt] (axis cs:0,\fmlinetop) -- (axis cs:0,\fmymax);
\node[anchor=base west, inner sep=0pt, font=\scriptsize, text=icInk, xshift=-2.5pt]
  at (axis cs:0,\fmylab) {Favors PsyPref\,$\rightarrow$};
% thin bars, thick bars, then markers: light first, blue on top; circle 3-35, square fixed 12
\foreach \c/\t/\s/\m in {\fmlight/in/w/-, \fmlight/in/f/+, icBlue/out/w/-, icBlue/out/f/+}{%
  \edef\fmplot{\noexpand\addplot[fm bar={\c}{0.5pt}, forget plot] table[col sep=comma, x=\t_\s,
    y expr=\noexpand\thisrow{y}\m\fmoff,
    x error plus expr=\noexpand\thisrow{\t_jhi_\s}-\noexpand\thisrow{\t_\s},
    x error minus expr=\noexpand\thisrow{\t_\s}-\noexpand\thisrow{\t_jlo_\s}]
    {\icorrtikzdata fig4m_rows.csv};}\fmplot
  \edef\fmplot{\noexpand\addplot[fm bar={\c}{1pt}, forget plot] table[col sep=comma, x=\t_\s,
    y expr=\noexpand\thisrow{y}\m\fmoff,
    x error plus expr=\noexpand\thisrow{\t_hi_\s}-\noexpand\thisrow{\t_\s},
    x error minus expr=\noexpand\thisrow{\t_\s}-\noexpand\thisrow{\t_lo_\s}]
    {\icorrtikzdata fig4m_rows.csv};}\fmplot
  \edef\fmplot{\noexpand\addplot[fm marks=\c, fm mark \s, forget plot] table[col sep=comma, x=\t_\s,
    y expr=\noexpand\thisrow{y}\m\fmoff] {\icorrtikzdata fig4m_rows.csv};}\fmplot}
% both methods at the cap: an open circle, nothing measured
\addplot[only marks, fm mark w, mark options={fill=white, draw=icBlue, line width=0.6pt}, forget plot]
  table[col sep=comma, x=cap_w, y expr=\thisrow{y}-\fmoff] {\icorrtikzdata fig4m_rows.csv};
\end{axis}
% labels of the blue cells: the value left of the thin bar, "both at 95 %" right of its circle, centred on their
% own marker's line, shifted by the script's vdy (pt, up positive) where something drawn is close
\def\fmvalue#1#2#3{% row, setting tag, marker offset (row units)
  \pgfmathsetmacro\fmvy{\fmB+(\fmymax-\icv{fmr}{#1}{y}-(#3))*\fmP+\icv{fmr}{#1}{vdy_#2}}%
  \ifnum\icv{fmr}{#1}{pos_#2}=3
    \node[inner sep=0pt, font=\scriptsize, text=icBlue, anchor=east] at (\icv{fmr}{#1}{vx_#2},\fmvy)
      {\ifnum\icv{fmr}{#1}{neg_#2}=1\textminus\fi\icv{fmr}{#1}{lab_#2}};%
  \fi
  \ifnum\icv{fmr}{#1}{pos_#2}=4
    \node[inner sep=0pt, font=\scriptsize, text=icBlue, anchor=west] at (\icv{fmr}{#1}{vx_#2},\fmvy) {both at 95\,\%};%
  \fi}%
% group headers and row labels, left of the axis
\foreach \i in {0,...,\fmN}{%
  \pgfmathsetmacro\fmy{\fmB+(\fmymax-\icv{fmr}{\i}{y})*\fmP}%
  \ifnum\icv{fmr}{\i}{is_header}=1
    \node[anchor=mid west, inner sep=0pt, font=\scriptsize\bfseries, text=icInk]
      at (0,\fmy) {\fmlab{\icv{fmr}{\i}{key}}};%
  \else
    \ifnum\icv{fmr}{\i}{bold}=1 \def\fmfont{\scriptsize\bfseries}\def\fmx{0}%
    \else \def\fmfont{\scriptsize}\def\fmx{5}\fi
    \node[anchor=mid west, inner sep=0pt, font=\fmfont, text=icInk]
      at (\fmx,\fmy) {\fmlab{\icv{fmr}{\i}{key}}%
        \ifnum\icv{fmr}{\i}{n_merged}>1\ (\icv{fmr}{\i}{n_merged})\fi
        \ifnum\icv{fmr}{\i}{tie}=1\,\textdagger\fi};%
    \fmvalue{\i}{w}{-\fmoff}%
    \fmvalue{\i}{f}{\fmoff}%
  \fi}
% key (shape = the setting's peak torque), on the line of the zero line's label, left
\pgfmathsetmacro\fmyk{\fmB+(\fmymax-\fmylab)*\fmP}%
\begin{scope}[font=\scriptsize, every node/.style={anchor=base west, inner sep=0pt, text=icInk}]
  \node (fmk1) at (0,\fmyk) {Peak torque};
  \fill[\fmlight] ([xshift=4.4pt, yshift=1.7pt]fmk1.base east) circle[radius=\fmmr pt];
  \node (fmk2) at ([xshift=7.3pt]fmk1.base east) {3--35\,Nm};
  \fill[\fmlight] ([xshift=4.4pt-\fmms pt, yshift=1.7pt-\fmms pt]fmk2.base east)
    rectangle ([xshift=4.4pt+\fmms pt, yshift=1.7pt+\fmms pt]fmk2.base east);
  \node at ([xshift=7.3pt]fmk2.base east) {Fixed 12\,Nm};
\end{scope}
% x label, one line, centred under the zero line
\pgfmathsetmacro\fmzero{\fmL-\fmxmin*\fmscale}%
\node[anchor=base, inner sep=0pt, font=\footnotesize, text=icInk] at (\fmzero,1.9)
  {$\Delta P_\mathrm{g}$ (percentage points)};
\pgfresetboundingbox
\useasboundingbox (0,0) rectangle (252,\fmB+\fmH);
\end{tikzpicture}%
